\documentclass[border=3pt]{standalone}
\usepackage{fontspec}
\setmainfont{Times New Roman}
\usepackage{pgfplots}
\pgfplotsset{compat=1.18}
\usetikzlibrary{arrows.meta,calc,positioning,shapes.geometric}
\definecolor{forest}{HTML}{346E57}
\definecolor{sage}{HTML}{E8EFEA}
\definecolor{clay}{HTML}{B77E60}
\definecolor{gray}{HTML}{718078}
\definecolor{rule}{HTML}{B7C3BA}
\definecolor{ink}{HTML}{000000}
\begin{document}
\hyphenpenalty=10000\exhyphenpenalty=10000
\fontsize{10}{12}\selectfont
\renewcommand{\small}{\fontsize{9}{11}\selectfont}
\renewcommand{\scriptsize}{\fontsize{8}{10}\selectfont}
\resizebox{16cm}{!}{%
\begin{tikzpicture}[x=5cm,y=2.7cm,axis/.style={-{Latex[length=2.2mm]},draw=rule,line width=.7pt}]
\draw[axis] (-1.05,0)--(1.08,0);
\draw[axis] (0,-1.05)--(0,1.08);
\node[anchor=east,text width=2.4cm,align=center] at (-1.1,0) {Penawaran umum};
\node[anchor=west,text width=2.4cm,align=center] at (1.13,0) {Penanganan khusus};
\node[anchor=south,text width=2.4cm,align=center] at (0,1.13) {Koordinasi terpadu};
\node[anchor=north,text width=2.4cm,align=center] at (0,-1.13) {Pengaturan terpisah};
\fill[forest] (-0.65,-0.45) circle[radius=1.4pt];
\node[align=center,text width=2.6cm,anchor=south] at (-0.65,-0.33) {Alternatif A};
\fill[forest] (0.6,-0.5) circle[radius=1.4pt];
\node[align=center,text width=2.6cm,anchor=south] at (0.6,-0.38) {Alternatif B};
\fill[forest] (0.65,0.65) circle[radius=1.4pt];
\node[align=center,text width=2.6cm,anchor=south] at (0.65,0.77) {Usulan layanan};
\fill[forest] (-0.55,0.45) circle[radius=1.4pt];
\node[align=center,text width=2.6cm,anchor=south] at (-0.55,0.5700000000000001) {Alternatif C};
\end{tikzpicture}
}

\end{document}
